\section{Finite Cauchy--Schwarz equality}

\begin{theorem}[Equality in finite Cauchy--Schwarz]
    \label{thm:finite_cauchy_schwarz_equality}
    \lean{Finset.sq_sum_eq_card_mul_sum_sq_iff}
    \leanok
    Let $S$ be a finite set and let $(x_i)_{i\in S}$ be real numbers. Then
    \[
        \left(\sum_{i\in S}x_i\right)^2
        =|S|\sum_{i\in S}x_i^2
    \]
    if and only if $x_i=x_j$ for every $i,j\in S$.
\end{theorem}

\begin{proof}\leanok
    The variance identity gives
    \[
        \sum_{i,j\in S}(x_i-x_j)^2
        =2\left(
          |S|\sum_{i\in S}x_i^2-
          \left(\sum_{i\in S}x_i\right)^2
        \right).
    \]
    Since every summand on the left is non-negative, equality in
    Cauchy--Schwarz holds precisely when every difference $x_i-x_j$ vanishes.
    This argument also includes the empty and singleton cases.
\end{proof}

This criterion is the equality condition invoked for the off-diagonal overlaps
and for the eigenvalues in \cite[Chapter~2, Proposition ``SIC POVMs'',
equations~(2.31)--(2.32)]{Wolf2012Quantum}.

\section{Trace-square and purity equality}

\begin{definition}[Rank-one orthogonal projection]
    \label{def:sic_rank_one_orthogonal_projection}
    \lean{IsRankOneOrthogonalProjection}
    \leanok
    \uses{def:orthogonal_projection_channel}
    A matrix $P\in\MN{D}$ is a rank-one orthogonal projection if
    $P=P^\dagger$, $P^2=P$, and $\rank P=1$.
\end{definition}

\begin{theorem}[Hermitian trace-square inequality]
    \label{thm:sic_hermitian_trace_square}
    \lean{Matrix.IsHermitian.trace_re_sq_le_card_mul_trace_sq_re}
    \leanok
    If $Q\in\MN{D}$ is Hermitian, then
    \begin{align}
        \tr(Q)^2\le D\tr(Q^2).
        \label{eq:sic_hermitian_trace_square}
    \end{align}
\end{theorem}

\begin{proof}\leanok
    \uses{thm:hermitian_trace_power_sum}
    Let $\lambda_0,\ldots,\lambda_{D-1}$ be the eigenvalues of $Q$.
    The spectral theorem and finite Cauchy--Schwarz give
    \begin{align}
        \tr(Q)^2
        =\left(\sum_i\lambda_i\right)^2
        \le D\sum_i\lambda_i^2
        =D\tr(Q^2).
    \end{align}
\end{proof}

\begin{theorem}[Equality in the Hermitian trace-square inequality]
    \label{thm:sic_hermitian_trace_square_equality}
    \lean{Matrix.IsHermitian.trace_re_sq_eq_card_mul_trace_sq_re_iff}
    \leanok
    A Hermitian matrix $Q\in\MN{D}$ satisfies equality
    in~(\ref{eq:sic_hermitian_trace_square}) if and only if
    $Q=r\Id$ for some $r\in\R$.
\end{theorem}

\begin{proof}\leanok
    \uses{thm:finite_cauchy_schwarz_equality, thm:hermitian_trace_power_sum}
    Equality holds precisely when all eigenvalues of $Q$ coincide. If their
    common value is $r$, unitary diagonalization gives
    $Q=U(r\Id)U^\dagger=r\Id$. The converse follows by direct substitution.
    When $D=0$, every matrix is the zero matrix, so the same conclusion holds.
\end{proof}

\begin{theorem}[Trace bound at unit purity]
    \label{thm:sic_psd_unit_purity_trace}
    \lean{Matrix.PosSemidef.one_le_trace_re_of_trace_sq_eq_one}
    \leanok
    If $P\in\MN{D}$ is positive semidefinite and $\tr(P^2)=1$, then
    $\tr(P)\ge 1$.
\end{theorem}

\begin{proof}\leanok
    \uses{thm:hermitian_trace_power_sum}
    Write the non-negative eigenvalues of $P$ as $\lambda_i$. Then
    \begin{align}
        1=\sum_i\lambda_i^2
        \le\left(\sum_i\lambda_i\right)^2
        =\tr(P)^2.
    \end{align}
    Since $\tr(P)=\sum_i\lambda_i\ge0$, the assertion follows.
\end{proof}

\begin{theorem}[Equality at unit purity]
    \label{thm:sic_psd_unit_purity_equality}
    \lean{Matrix.PosSemidef.trace_re_eq_one_iff_isRankOneOrthogonalProjection}
    \leanok
    \uses{def:sic_rank_one_orthogonal_projection}
    Suppose that $P\in\MN{D}$ is positive semidefinite and $\tr(P^2)=1$.
    Then $\tr(P)=1$ if and only if $P$ is a rank-one orthogonal projection.
\end{theorem}

\begin{proof}\leanok
    \uses{thm:hermitian_trace_power_sum}
    If $\tr(P)=1$, the non-negative eigenvalues satisfy
    \begin{align}
        \sum_i\lambda_i=\sum_i\lambda_i^2=1.
    \end{align}
    Each $\lambda_i$ lies in $[0,1]$, and
    $\sum_i\lambda_i(1-\lambda_i)=0$. Hence every eigenvalue is either zero
    or one. Thus $P^2=P$, and its rank equals its trace, namely one. Conversely,
    a rank-one orthogonal projection has one unit eigenvalue and all remaining
    eigenvalues zero, so its trace is one.
\end{proof}
